% Generated draft: requires mathematical and source review.
\documentclass[11pt]{article}
\usepackage[margin=25mm]{geometry}
\usepackage{fontspec}
\setmainfont{Latin Modern Roman}
\usepackage{amsmath,amssymb,mathtools,unicode-math}
\setmathfont{Latin Modern Math}
\usepackage{enumitem,needspace,xurl,hyperref}
\setlength{\parindent}{0pt}
\setlength{\parskip}{5pt}
\setlength{\emergencystretch}{3em}
\newcommand{\N}{\mathbb{N}}
\newcommand{\Z}{\mathbb{Z}}
\newcommand{\Q}{\mathbb{Q}}
\newcommand{\R}{\mathbb{R}}
\newcommand{\C}{\mathbb{C}}
\newcommand{\F}{\mathbb{F}}
\newcommand{\A}{\mathbb{A}}
\newcommand{\PP}{\mathbb{P}}
\newcommand{\E}{\mathbb{E}}
\newcommand{\eps}{\varepsilon}
\newcommand{\dd}{\,\mathrm{d}}
\begin{document}
\section*{UnsolvedMath ID: 30004159}
\textbf{Source problem code:} catalogue entry
\subsection{Definitions and mathematical statement}
A nonzero holomorphic abelian differential on a smooth connected complex curve \(C\) of genus \(g\) is a section \(0\ne\omega\in H^0(C,\Omega_C^1)\). Its zero divisor has degree \(2g-2\). A zero signature \(\mu=(m_1,\ldots,m_n)\) records the multiplicities of its zeros, so that \(\sum_{i=1}^n m_i=2g-2\). Entries recording actual zeros are positive; allowing entries equal to zero instead records additional marked regular points. The source does not specify which convention or range of signatures is intended.

Under the conventional interpretation over \(\C\) with ordered marked points, the projectivized stratum \(\mathbb P\Omega\mathcal M_g(\mu)\) parametrizes isomorphism classes of tuples
\[
(C,p_1,\ldots,p_n,[\omega]),\qquad
[\omega]\in\bigl(H^0(C,\Omega_C^1)\setminus\{0\}\bigr)/\C^*,
\qquad
\operatorname{div}(\omega)=\sum_{i=1}^n m_i p_i,
\]
where \(C\) is smooth, connected and projective of genus \(g\), and the \(p_i\) are pairwise distinct. Isomorphisms preserve the marked points and the differential up to a nonzero scalar. This describes a conventional marked moduli problem; the supplied statement does not specify the base field, ordered versus unordered zeros, or whether the stratum denotes a stack or its coarse moduli space.

\textbf{Requested construction.} For the genus \(g\) and signature \(\mu\) intended in the source, construct a compactification \(\overline{\mathcal H}_{g,\mu}\) together with an open immersion
\[
j:\mathbb P\Omega\mathcal M_g(\mu)\hookrightarrow\overline{\mathcal H}_{g,\mu}
\]
whose image is dense, such that \(\overline{\mathcal H}_{g,\mu}\) is proper and smooth over the intended base and has a natural modular interpretation. In the conventional complex interpretation, properness and smoothness are relative to \(\operatorname{Spec}\C\). A modular interpretation should specify the geometric objects and their families, including boundary objects, whose moduli recover the compactification and restrict to the abelian-differential moduli problem on the open stratum.

The intended range of \((g,\mu)\), the ambient category in which properness and smoothness are required, and the precise meaning of a natural modular interpretation remain unspecified. Consequently, the displayed construction is a conditional formulation of the supplied target, rather than a claim that these choices have been resolved.
\subsection{Short English statement}
Construct a proper, smooth compactification of the projectivized stratum of holomorphic abelian differentials with prescribed zero multiplicities, and describe its points and families by a natural moduli problem extending the interior one. The source leaves the parameter range, base, marking convention, category of moduli objects, and required boundary interpretation unresolved. Its \texttt{partially\_solved} label is supplied dated curation, not an independently verified mathematical conclusion.
\subsection{Sources}
\begin{itemize}
\item Supplied local UnsolvedMath registry record, permanent ID \(30004159\), title \emph{Smooth Modular Compactifications of Abelian-Differential Strata}, in \emph{Oberwolfach Reports Open Problems}; exact source problem code \texttt{catalogue entry}. The record gives proposed year \(2019\) and names no individual proposer.
\item Supplied primary-report citation: \emph{Logarithmic Enumerative Geometry and Mirror Symmetry} (2020), \emph{Oberwolfach Reports}, pp.\ 1639--1695, DOI \texttt{10.4171/owr/2019/27}. Supplied source URL and literature-source label ``OWR primary report'': \url{https://doi.org/10.4171/owr/2019/27}. The record describes these reports as documenting workshops and their open-problem discussions.
\item Supplied literature-source label ``moduli space of multi-scale differentials'': \url{https://arxiv.org/abs/1910.13492}. No authors or further bibliographic details are supplied in the record.
\item Literature assessment dated exactly \texttt{2026-08-22}: the supplied curation labels the problem partially solved and reports constructions of logarithmic and related compactifications, while stating that simultaneous properness, smoothness, and a natural modular interpretation are not established in the extracted generality. This is supplied dated curation, not independently verified literature assessment; the record explicitly requests independent verification of time-sensitive status and cited claims.
\item The supplied statement status is \texttt{corrected\_verified}; its accompanying explanation says that notation, typography, and sentence structure were repaired while mathematical content was preserved. This registry label does not establish mathematical correctness or independently verify the cited sources.
\end{itemize}
\end{document}
